\section{Good paths give the required segment ranks}
\label{sec:good-paths}

\status{The actual \texttt{RestMem}, \texttt{RestInd},
\texttt{GoodPathWeight}, \texttt{BetaRatioBound} and \texttt{RestSeg}
are proved under \texttt{Dat.Hyp} and $C_{\mathcal L},C_k\geq1088$.
The proof uses the earlier IP integral path basis, CB dimension formula,
SD6 residue count and the RS2--RS3 arithmetic. Supplying the global
frame and parameter data remains a separate step.}

\paragraph{GP1--GP2: the actual paths and their total mass.}
Assume \texttt{Dat.Hyp} and $C_{\mathcal L},C_k\geq1088$.
In one segment write $x_0<\cdots<x_{N-1}$, $n_0=N-r$, $k=r+K$.
The first $r$ letters are the old frame; the last $r$ are the current
frame $\tau_t$.  Put $P_i=P_z(j,x_i)$, $Q=Q_{c,j}$, $m=kN_s+1$, and
\[
 B_i=\begin{cases}(B_{w,i}-m)_+,&i<r,\\b(x_i),&r\leq i<n_0,\end{cases}
 \qquad \beta=\sum_{i<n_0}B_i,\qquad
 R=\prod_{i<n_0}\binom{B_i+k-1}{k-1}.
 \tag{GP1}
\]
For counting, $b$ need only satisfy the source's intermediate-degree
condition \texttt{AdmLL}; membership below additionally requires
\texttt{BOK}.
The RS bounds give $k\leq B_i\leq2a/P_i$ and
$B_{i'}\geq(p^2/5)B_i$ for earlier letters $i<i'<n_0$ of the same
parity.  Also $N\leq N_s$, $N\geq3r$, and
$h:=K-2N\geq k/2$.  For $i\geq r$, let $o(i)$ be the $r$-th most
recent vector letter; it lies before $n_0$.

Take the literal signed horizontal-strip paths $\mu^0=0,\ldots,\mu^N$:
their first $n_0$ strip sizes are $B_i$, and their terminal conditions are
$\mu^{n_0+t}_l=\beta$ for $l<t$.  Let $\nu,\kappa$ be the first $r$
and last $K$ rows.  A path has weight $w(\mu)=d_K(\kappa^N)$, the
number of its standard normal arrays.  Signed Pieri at the fixed-size
steps and GT branching at the terminal steps give
\[
 \sum_{\mu}w(\mu)=R,\qquad
 \#\{\mu\}\leq\exp(2N^2L),\qquad L=\log(1+2aN).
 \tag{GP2}
\]
For the second bound, at each stage only the first $N$ and last $N$
coordinates can be nonzero; each has at most $1+\beta\leq1+2aN$
possible values.  Terminal branching is exact: deleting the forced
$\beta$ entries turns every terminal strip into one interlacing row.

\paragraph{GP3--GP6: why the bad paths have small mass.}
At stage $i$ delete its $(i-n_0)_+$ forced heads, obtaining $\eta^i$.
Let $T_+(i,u)$ be the sum of the positive parts after coordinate $u$,
and $T_-(i,u)$ the same sum for the reversed negative parts.
Let $S_\pm(i,u)$ be the suffix sums of the reverse-chronological lists
of earlier preterminal vector/covector degrees.
Horizontal-strip interlacing gives
\[
 T_\pm(i,u)\leq S_\pm(i,u).
 \tag{GP3}
\]
Indeed an active step prepends its degree to the reference list and
satisfies $T(i+1,u+1)\leq T(i,u)$ and
$T(i+1,0)\leq T(i,0)+B_i$; an opposite-sign or terminal step decreases
the relevant tails.  Thus a deficit $X$ at a suffix persists to the
endpoint, shifting its index once for each later active step.

Set $f_h(x)=\sum_{t=1}^h\log(1+x/t)$.  For a reference list
$b_u\geq h$, $b_u\geq4b_{u+1}$ and nonnegative entries $a_u$ satisfying the
suffix bounds, concavity and summation by parts give
\[
 \begin{split}
 \sum_u\bigl(f_h(a_u)-f_h(b_u)\bigr)
 &\leq\sum_u(g_u-g_{u-1})\Delta_u,\\
 g_u=f'_h(b_u),\quad g_{-1}&=0,\quad
 \Delta_u=\sum_{l\geq u}(a_l-b_l)\leq0.
 \end{split}
 \tag{GP4}
\]
Since $h/(x+h)\leq f'_h(x)\leq h/(x+1)$, the coefficients satisfy
$g_u-g_{u-1}\geq h/(4b_u)$.  A deficit $X$ therefore contributes
at most $-Xh/(4b_u)$.

The endpoint $\lambda=\kappa^N$ has at most $N$ positive and $N$
negative entries.  In its Weyl product, pairs of nonzero entries cost
at most $2N^2L$.  For each nonzero entry, its interactions with zero
entries are bounded by $f_h(|\lambda_l|)+2NL$: the distances beyond
$1,\ldots,h$ number at most $2N$.  Hence
\[
 \begin{aligned}
 \log w(\mu)&\leq
 \sum_l f_h((\lambda_l)_+)+\sum_l f_h((-\lambda_l)_+)+6N^2L,\\
 \log R&\geq\sum_{i<n_0}f_h(B_i).
 \end{aligned}
 \tag{GP5}
\]

Put $\Xi_i=\mathcal L a/(pkP_i)$.  Call a path typical if a covector
step loses at most $\Xi_i$ positive mass, a vector step $i\geq r$
loses at most $\Xi_i$ negative mass, and
$\mu^i_{r-1}\geq2B_{o(i)}/3$.  Failure of a mass condition creates
a suffix deficit $\Xi_i$.  Its reference entry is an earlier
$B_l\leq4a/(pP_i)$, so (GP4) costs at least $\mathcal L h/(16k)$.
Failure of the last condition creates a deficit $B_{o(i)}/3$ at that
reference entry, costing $h/12$.  Combining (GP2)--(GP5),
\[
 \sum_{\text{atypical }\mu}w(\mu)
 \leq R\exp\!\left(8N^2L-\frac{\min(\mathcal L,k)}{32}\right)
 \leq R a^{-N_s^2}.
 \tag{GP6}
\]
Here $L\leq2\log a$ and
$\min(\mathcal L,k)\geq1088N_s^2\log a$; the estimate is stronger
than the source's exponent with $11N^2L$.

\paragraph{GP7--GP9: good paths yield actual independent generators.}
Tracing row $r-1$ backwards through its $r-1$ later vector steps and
summing the geometric degree separation gives
$\mu^i_{r-1}\leq(1+5/(p^2-5))B_{o(i)}<2B_{o(i)}$.
Interlacing then bounds the normal-degree increase at a vector step
by $\Xi_i+2B_{o(i)}$, and the plane-degree decrease at a covector
step by $\Xi_i$.  These are precisely the source's
$\mathrm{thc}=2B_{o(i)}/3$ and
$\mathrm{oc}_i=\Xi_i+\mathbf1_{x_i\text{ even}}4a/P_{o(i)}$
bounds.  Thus the good-path/normal-array set $T$ satisfies
$|T|\geq(1-a^{-N_s^2})R$.

For $(\mu,I)\in T$, take the actual auxiliary exponents
$u_{tc}+2\beta<Q$ with terminal row degrees
$d_{n_0+t}+\sum_cu_{tc}=b(\tau_t)$.  The generators are exactly
\[
 \begin{aligned}
 F_{\mu,I,u}&=C_b\prod_{t,c}X_{\tau_t,c}^{u_{tc}}
 \Psi_{\rm pl}(\nu,I_{\rm can})\Psi_{\rm nor}(\kappa,I),\\
 C_b&=\det(X_{W})^m\prod_{x\in W}X_{x,0}^{\,b_x-B_{w,x}}.
 \end{aligned}
 \tag{GP7}
\]
The IP prefix bounds and the crossing inequalities
$\kappa^{i+1}_0\leq\nu^i_{r-1}$ for vectors,
$\kappa^i_0\leq\nu^{i+1}_{r-1}$ for covectors prove the literal
$XX'$ conditions: $\det(X_W)^m$ supplies the additional $m$.
The normal polynomial has no old-frame variables.  Row degrees add
to $b$, and SD6 gives $d_{n_0+t}\leq2\beta$, so every terminal plane
exponent is $<Q$.  Hence $F_{\mu,I,u}\in M_j(b)$.

For independence, retain only the antidiagonal variable
$Y_t=X_{\tau_t,r-1-t}$ in each terminal plane row of the path polynomial.
Column expansion makes it a signed monomial in the $Y_t$ times a
truncated-path standard polynomial.  Its retained column sets record
the time of every terminal column rise, hence recover the full path.
The IP integral basis and disjoint terminal variables prove independence
of these specialized products, including the auxiliary monomials.
In a relation before specialization, scale the killed variables by $T$
and take the least auxiliary $T$-degree with nonzero coefficient;
this gives a forbidden relation among the specialized products.
Finally cancel the common nonzero $C_b$.  This proves
\texttt{RestMem} and \texttt{RestInd} over the actual DVR.

\paragraph{GP10--GP12: the residue bound and \texttt{RestSeg}.}
Let $D=\operatorname{dist}(j,\tau_0)$, $g=z_{j+1}-z_j$.
The preceding gap and distinct distances give
$\beta\leq4a/p^{R_j+D+L_1+1}$.  RS3 gives
$Q=p^{g+R_{j+1}-R_j}>a/(8pYk^2N_s p^{D+R_j})$.  Therefore
\[
 \frac\beta Q<\frac{32Yk^2N_s}{p^{L_1}}
 \leq\frac1{4pY}\leq\frac1{4Y}.
 \tag{GP8}
\]
The elementary SD6 residue count now supplies, for every terminal row
and every desired residue, at least
$Q^{r-1}\max(0,1-r/(2Y))$ choices of $u$.
Their resulting degrees $e_t=d_{n_0+t}+\sum_cu_{tc}$ satisfy
$0\leq e_t<rQ$.  Set $G_j(b)=|\mathrm{Idx}(j,\mathrm{restr}\,b)|$.
This index set equals $\mathrm{Idx}(j,b)$: old-frame inputs use
$B_w-m$, independently of the old values of $b$.
Membership, independence and finite-dimensionality give
$G_j(b)\leq\operatorname{rank}_O M_j(b)$ for every \texttt{BOK} $b$;
the count gives
\[
 \sum_{\substack{0\leq e_t<rQ\\e_t\equiv\rho_t\pmod Q}}
 G_j(b[e])
 \geq(1-a^{-N_s^2})R
       \left[Q^{r-1}\max\!\left(0,1-\frac r{2Y}\right)\right]^r.
 \tag{GP9}
\]
The function depends only on late row degrees.  These are exactly
the clauses of \texttt{RestSeg 1088}; no path-abundance or numerical
ratio hypothesis remains unproved.
